% !TeX root = ../nc-cubic-non-as.tex
\subsection{零係数の台の有限分類と多零係数型の排除}
\label{subsec:ou-zero-supports}

$c=(c_{ijk})$を三ブロック箙の関係テンソルとし、
$z(c)=\#\{(i,j,k):c_{ijk}=0\}$とする。
以下では標数零の代数閉体を考える。
独立な三添字の置換と添字方向の置換による台の同値を用いる。
偶置換の方向変更は巡回foundationの変更に対応し、
奇置換はこれに経路の反転（反対代数）を加えたものである。
従って下記の三次syzygyの必要条件と固定端点の経路商の
次元超過を判定する上では、いずれも許してよい。
内在的条件のheartの同値を反対代数から直接導くわけではない。

\begin{proposition}[台の必要条件の有限列挙]
\label{prop:ou-zero-support-enumeration}
very strong導来helixを持つなら$z(c)\leq12$である。
さらに補題\ref{lem:ou-cyclic-cubic-syzygy}と
系\ref{cor:ou-cyclic-support-obstruction}の
台だけに依存する必要条件を満たす台の数とその同値類数は
次の通りである。数は関係代数や曲面圏の同型類数ではない。
\[
\begin{array}{r|r|r}
 z & \text{ラベル付き台} & \text{台の同値類}\\ \hline
0&1&1\\
1&27&1\\
2&351&3\\
3&2898&9\\
4&16659&30\\
5&68256&80\\
6&196110&201\\
7&374490&344\\
8&428004&387\\
9&242688&217\\
10&49680&53\\
11&1944&4\\
12&72&1
\end{array}
\]
\end{proposition}
\begin{proof}
各方向の固定添字に対応する$3\times3$のsliceについて、
その六頂点二部graphは連結でなければならない。
よって各sliceは少なくとも五辺を持ち、一方向の三sliceを
合計すれば非零係数は15個以上、すなわち$z(c)\leq12$を得る。
異なる二sliceの台を$a,b$と書くと、三次syzygyの式
$\alpha_k c_{ijk}=\beta_j c_{i'jk}$において、
共通辺は二変数を非零倍で結び、一方だけにある辺は
対応する一変数を零に強制する。
強制零は共通辺の連結成分全体に伝播する。
強制零を含まない木成分（孤立頂点も含む）は、
どの非零係数に対しても一自由度の解を与える。
従ってこのような成分がある台は排除される。
両方のsliceの順序を調べる必要がある。
cycleを含む成分で係数の比が整合するか否かは
この台だけの判定には含めない。

再現コード\path{computations/ou_zero_supports.cpp}は
$3\times3$の全512台から連結な205台を抽出し、
その三つの順序付き組から一方向のsliceを構成する。
残る二方向の全sliceの連結性と、全方向で異なるsliceの
両順序に対する上の条件を調べる。
その後$(\mathfrak S_3)^3\rtimes\mathfrak S_3$の全1296元で
台を置換して同値類を数える。
台は27ビットの整数として表現し、ビット$9i+3j+k$が
1であることを$c_{ijk}\ne0$とする。
この有限列挙によって表を得る。
\end{proof}

\begin{theorem}[九個以上の零係数を持つvery strong型の制約]
\label{thm:ou-large-zero-path-exclusion}
三ブロック箙の導来helixがvery strongなら$z(c)\leq9$である。
$z(c)=9$なら独立な三添字の置換により
\[
 c_{ijk}=0\quad\Longleftrightarrow\quad i+j+k=0
 \quad\text{in }\mathbb Z/3\mathbb Z
\]
となる。この台では非零18係数のうち17係数を
矢印の非零スカラー倍で1に正規化できる。
\end{theorem}
\begin{proof}
命題\ref{prop:ou-zero-support-enumeration}の各台に対して
矢印のスカラー倍の作用をまず用いる。
非零係数の集合を$T$とし、指数行列$B_T$を
$|T|\times27$の整数行列として定める。
各行$(i,j,k)\in T$は三矢印$(i,j),(j,k),(k,i)$に
対応する成分が1、他が0である。
対応するtorusの作用はこの指数行列を持つmonomial mapである。
$\operatorname{rank}_{\mathbb Q}B_T=r_T$なら
独立な$r_T$行に対応する非零係数を同時に1にできる。
実際、この部分行列のSmith標準形により、
標数零の代数閉体上では非零整数冪写像が全射であることから
対応するtorusへの射が全射となる。
残る$|T|-r_T$個の非零係数を独立な変数とする族は、
すべての係数の選択を正規化した後に含む。
これが同型類の一意なパラメータ表示であるとは主張しない。

正規化した各族の固定端点の長さ$m$の経路商を計算する。
前節と同様、内部の$m-1$添字を辞書式順序に並べ、
各内部位置とそれ以外の内部添字を固定して
その位置の三通りの選択に対応する二次関係を一行とする。
従って列数は$3^{m-1}$、行数は高々$(m-1)3^{m-2}$である。
行列は係数変数の多項式行列であり、
その有理関数体上のrankはすべての特殊化のrankの上界である。
従って経路商の次元の下界を得るために
generic rankを厳密に計算すれば十分である。
$z=9,10,11,12$について$m=3,4,5$の全巡回、
全九端点を検査し、最初の次元超過を用いる。
very strong性なら$m=3$では$3+\delta_{ij}$、
$m=4$では5、$m=5$では7である
（補題\ref{lem:ou-cyclic-cubic-syzygy}のEuler行列）。
有限計算の結果は次の通りである。
\[
\begin{array}{r|r|r|r|r}
 z & |T|-r_T & m=4\text{で排除} & m=5\text{で排除}
 & \text{残余}\\ \hline
12&0&1&0&0\\
11&0&4&0&0\\
10&0&52&1&0\\
9&0&201&10&0\\
9&1&5&0&1
\end{array}
\]
計算は\path{computations/verify_ou_zero_supports.py}で再現できる。
整数指数行列は有理数体上で行簡約し、
経路関係行列のrankは有理数体または一変数有理関数体上で
正確な線形代数により求める。
有限体への特殊化や数値近似はこの排除に用いない。
$m=3$の計算ではここで列挙された各族に新たな排除はない。
$m=4,5$の排除は全特殊化に対する次元下界による。

唯一の残余の台はビット整数108455662であり、
その非零成分は$i-j+k\ne0$に対応する。
$j$を$-j$と取り直せば主張の対称な台となる。
他の方向置換をした代表についても、この条件は対称なので、
独立な添字変更だけで同じ台にできる。
上のrank計算で$r_T=17$であるから、係数を17個正規化できる。
\end{proof}

\begin{remark}
九個の零係数を持つ残余族がvery strongであることは、
長さ5までの次元検査からは従わない。
この族の内在的三次条件については次の定理で別に排除する。
\end{remark}

\begin{theorem}[七・八零係数の台の経路次元による制約]
\label{thm:ou-seven-eight-zero-path-exclusion}
very strong導来helixを持つ三ブロック型について、
$z(c)=8$なら前定理のLatin square型の九つの零係数のうち
一つだけを非零にした台に限られる。この台の非零19係数中18個は
矢印の非零スカラー倍で1にできる。
$z(c)=7$なら、命題\ref{prop:ou-zero-support-enumeration}の
台の同値を除いて、次の八つの台に限られる：
\[
\begin{array}{r|c}
\text{台のビット整数}&\text{正規化後の自由係数数}\\ \hline
116309742&1\\
116834030&1\\
117368558&1\\
125234926&2\\
125531767&1\\
128957047&1\\
129972855&1\\
133118583&1
\end{array}
\]
この列挙は実現可能性を主張しない。
\end{theorem}
\begin{proof}
前定理と同じ指数行列を用いた正規化を、
$z=7,8$の全台に対して行う。
$\mathbb Q$または自由係数の有理関数体上で、
長さ$3,4,5,6,7$の全巡回・全九端点の経路商を調べる。
長さ6の必要次元は$9+\delta_{ij}$、長さ7では12である。
これは補題\ref{lem:ou-cyclic-cubic-syzygy}の$D,S_0$から
$(S_0^2)^{\mathsf t}D$の対応するブロックを計算して得る
（定理\ref{thm:ou-dense-no-triple-grid-failure}の証明参照）。
最初の次元超過によって排除する台の数は
\[
\begin{array}{r|r|r|r|r|r|r}
z&\text{台数}&m=4&m=5&m=6&m=7&\text{残余}\\ \hline
8&387&319&55&10&2&1\\
7&344&192&103&26&15&8
\end{array}
\]
となる。長さ3では新しい排除はない。
\path{computations/verify_ou_zero_supports.py}が
この表と残余のビット整数を再現する。
計算の高速化のため自由係数を2,3,\ldotsとする有理数特殊化も用いる。
経路関係行列は多項式行列だから、特殊化のrankはgeneric rank以下である。
従ってある固定端点について特殊化で商次元が必要値以下なら
generic商次元の超過はあり得ない。
逆に特殊化で商次元が超過した場合は、
元の有理関数体上のrankを厳密計算し、超過が残る場合だけ排除する。
従って排除の全ては有理関数体（または有理数体）上のrank上界によるもので、
一つの数値例から一般族を排除するものではない。
残余とはこれらの長さの経路次元による全族排除ができない台である。

$z=8$の唯一の残余は108979950であり、非零係数の集合は
ビット整数108455662（$i-j+k\ne0$）の非零係数に
$(2,0,1)$を加えたもの。指数行列のrankは18である。
$j$の符号を反転するとLatin square型の零を一つ除いた台となる。
$z=7$の残余とその自由係数数も同じ計算から得る。
\end{proof}

\begin{theorem}[対角零・一方向巡回零の台の内在的排除]
\label{thm:ou-latin-square-periodic-exclusion}
添字を$\mathbb Z/3\mathbb Z$とし、
独立な三添字の置換の後に次を満たす三ブロック関係代数を考える：
\[
 c_{iii}=0,\qquad c_{ijk}\ne0\quad(\#\{i,j,k\}=2),
 \qquad c_{i,i+\eta,i+2\eta}=0\quad(i\in\mathbb Z/3\mathbb Z)
\]
となる$\eta\in\{1,-1\}$が存在する。
残る三つの相異なる添字の係数は任意（零でもよい）。
すべての非零係数の値も任意とする。
標準noetherian OU helix比較を持つ曲面圏$\mathcal C$が
対応するなら、$\mathcal C$は$\mathfrak I_3$に属さない。
特にLatin square型の九零および八零の唯一の残余を含む。
\end{theorem}
\begin{proof}
$i=k$の関係は$j=i$の係数だけが零であり、
$j\ne i$の二経路の非零係数による二項関係である。
$i\ne k$では$j=i,k$に対応する二係数は非零であり、
第三の係数は任意でよい。
従って九頂点の反変表現$Q_i$を、
頂点$U_i,V_i,W_i$で一次元、他で零とし、
対角二矢印の作用を恒等とすることで定義できる。

補題\ref{lem:ou-iv-periodic-three}の証明をこの場合に適用する。
具体的に射影分解は
\[
 0\longrightarrow P_{U_i}
 \longrightarrow \bigoplus_{j\ne i}P_{V_j}
 \longrightarrow P_{W_i}\longrightarrow Q_i\longrightarrow0
\]
であり、左の射の二成分は$A_{ij}$の
$c_{iji}$倍、中央の射は$B_{ji}$である。
$U_l$で$l\ne i$なら、経由頂点$j=i$の係数が非零なので
他の二経路は二次元商空間の基底になる。
$l=i$では他の二経路の唯一の関係が左の射の像である。
この頂点別の検証により完全性を得る。
Serre関手を施した入射複体についても同様に、
$W_l$で$l\ne i$なら経由頂点$j=i$の係数が非零で、
$l=i$では唯一の対角関係で完全性を得る。
従って$sQ_i\simeq Q_i$である。
これは係数の符号や比によらない。
補題\ref{lem:ou-iv-periodic-simple}の証明は
この周期性、九射影に関する次元ベクトルとEuler行列だけを
用いているので、標準比較の下で$Q_i$は
heartの互いに非同型な単純対象となる。

任意長の有限長商を次のように構成する。
$D_m=k[t]/(t^m)$とし、全九頂点に$D_m$を置く。
各矢印を次のスカラーの乗法として定める（右加群の規約）。
\[
 A_{ij}=
 \begin{cases}
 1&i=j,\\ t&j=i+\eta,\\0&\text{otherwise},
 \end{cases}
 \qquad
 B_{jk}=
 \begin{cases}
 1&j=k,\\
 -c_{jkk}/c_{jjk}\,t&k=j+\eta,\\
 0&\text{otherwise}.
 \end{cases}
\]
添字はすべてmodulo 3である。
$i=k$では非対角の二経路はともに零となる。
実際、一方向の巡回矢印のみを非零にしたので、
二つの非対角矢印を続けて元の添字に戻ることはない。
$i\ne k$の関係の$j=i,k$から来る二項は、
$k=i+\eta$なら定義により打ち消し合い、それ以外はともに零。
三添字が相異なる項で二矢印の積が非零となり得るのは
$(i,j,k)=(i,i+\eta,i+2\eta)$の場合だけだが、
その係数は仮定により零である。
従って全ての関係式は満たされ、この表現$T_m$は
任意の$m$で定義できる。

$t$進filtrationの各逐次商は$Q=Q_0\oplus Q_1\oplus Q_2$である。
従ってheartは拡大で閉じていることから$T_m$はheartに属し、
組成因子は各$Q_i$が$m$個である。
$\operatorname{Hom}(T_m,Q_i)$は一次元である。
実際、$U_i$の値への線型写像は、非対角矢印
$A_{i,i+\eta}$が$t$を作用させることから$tD_m$を消す。
対角矢印によって$V_i,W_i$での写像も定まり、
剰余写像$T_m\to Q\to Q_i$がHomを生成する。

$F=\bigoplus_i P_{W_i}$とおけば$F\twoheadrightarrow Q$があり、
$\operatorname{Ext}^1(F,Q_i)=0$。
有限filtrationを用いて$F\to Q$を$F\to T_m$へ持ち上げられる。
もしこれが全射でなければ、その余核の単純商$Q_i$への射は
$T_m\to Q$を経由するので$F\to Q$の全射性に反する。
従って固定した$F$から任意の$T_m$への全射を得る。

ここで定理\ref{thm:ou-iv-periodic-exclusion}の
Euler形式と内在的三次条件による障害を適用する。
同定された数値的rank$r$は$\sum_i\chi(-,Q_i)$であり、
仮定された内在的標識$\mathcal O$について$r([\mathcal O])=1$。
同定理の
\eqref{eq:ou-iv-periodic-ext-self-bound}と
\eqref{eq:ou-iv-periodic-ext1-exact}により、
各固定heart対象への$\operatorname{Ext}^1(s^{-n}\mathcal O,-)$の
次元は$n\gg0$で有界であり、固定完全列に対し漸近的に完全。
周期性と$F\twoheadrightarrow T_m$から
\[
 m\sum_i\dim\operatorname{Ext}^1(\mathcal O,Q_i)\leq C_F
\]
を得る。$C_F$は$m$によらず、従って全
$\operatorname{Ext}^1(\mathcal O,Q_i)=0$。
同定理の負次数Hom消滅・Serre双対性から
$\operatorname{Ext}^2(\mathcal O,Q_i)=0$、
AZ1と周期性から$\operatorname{Hom}(\mathcal O,Q_i)\ne0$。
よって
\[
 1=r([\mathcal O])=\sum_i\chi(\mathcal O,Q_i)\geq3
\]
となって矛盾する。
この議論は内在的標識のexceptionalityを追加仮定しない。
\end{proof}

\begin{corollary}[内在的三次型の零係数上界]
\label{cor:ou-intrinsic-at-most-seven-zeros}
標準noetherian OU helix比較を持つ内在的三次型の
very strongなAOU三ブロック箙型では、$z(c)\leq7$である。
$z(c)=7$の場合、定理\ref{thm:ou-seven-eight-zero-path-exclusion}の
八つの残余のうち125234926の台も内在的三次条件から排除される。
\end{corollary}
\begin{proof}
定理\ref{thm:ou-large-zero-path-exclusion}で10個以上を排除し、
9個の残余および定理\ref{thm:ou-seven-eight-zero-path-exclusion}の
8個の残余を定理\ref{thm:ou-latin-square-periodic-exclusion}で排除する。
後者はその台の三添字の符号・平行移動を取り直して、
零の対角三成分と一方向の巡回三成分を持つ形にできるためである。
7零の125234926についても同じ台の条件が成立する。
このことは独立な添字置換の全$6^3$通りを列挙すれば検証できる。
\path{computations/verify_ou_zero_supports.py}は残余台について
この条件の判定を行う。条件は添字方向の置換でも保存される
（巡回の向きは逆転してよい）ので、反対代数への内在的条件の
移行を仮定する必要はない。
\end{proof}

\subsection{七零係数の二つの一係数族の認識}
\label{subsec:ou-seven-zero-curves}

\begin{proposition}[七零係数の滑らかな切断と節点切断]
\label{prop:ou-seven-zero-curve-recognition}
定理\ref{thm:ou-seven-eight-zero-path-exclusion}の台116309742,
116834030を考える。非零係数は指数行列の独立行を選んで
19個を1に正規化でき、残る一係数を$u\in k^\times$とする。
具体的に、台116309742では$u=c_{211}$、台116834030では
$u=c_{222}$として他の非零係数を1に取れる。
前者の$u\ne1$、後者の全ての$u\ne0$において
既約平面三次曲線による九線束表示を持つ。
従ってvery strong導来helix・標準noetherian OU比較・内在的三次条件
の下で標準AS四次元正規三次商として認識される。
\end{proposition}
\begin{proof}
定理\ref{thm:ou-seven-eight-zero-path-exclusion}と同じ行順序での
指数行列の行簡約により、自由係数の位置は上記の通りである。
係数行列の六座標を命題\ref{prop:ou-full-support-determinant}の
$(z_0,z_1,z_2,z_3,z_4,z_5)=(ad,bc,a,b,c,d)$の順に取る。
台116309742では三線形式の係数行列は
\[
 L=\begin{pmatrix}
 0&0&-1&1&1&0\\
 1&0&-1&0&u&0\\
 0&-1&0&0&0&-1
 \end{pmatrix}.
\]
最初の三列は独立であり、零平面を
\[
 (z_0,z_1,z_2,z_3,z_4,z_5)
 =(x+(1-u)y,-z,x+y,x,y,z)
\]
で表示できる。三次切断の式は
\[
 f=x^2y+(1-u)xy^2+(x+y)z^2=0 .
\]
偏微分$f_z=2(x+y)z$を用いる。
$z\ne0$の特異点なら$x+y=0$、
$f=0$から$ux^3=0$、従って$x=y=0$となるが、
その点で$f_x=z^2\ne0$である。
$z=0$では二変数三次式$xy(x+(1-u)y)$の重根が
特異点となるため、$u\ne1$なら滑らかである。
$u=1$の点はここでは認識しない。

台116834030では
\[
 L=\begin{pmatrix}
 0&0&-1&1&1&0\\
 1&0&0&0&1&0\\
 0&-1&0&1&0&-u
 \end{pmatrix},
 \qquad
 (z_0,\ldots,z_5)=(-y,x-uz,x+y,x,y,z).
\]
やはり最初の三列が独立であり、三次切断は
\[
 f=u(x+y)z^2-x^2z-xy^2-xyz=0 .
\]
$x=0$上の点は$[0:1:0]$または$[0:0:1]$であり、
いずれも滑らかである。
$x=1$のchartにおいて$f_z=(1+y)(2uz-1)$。
特異点が$y=-1$を満たすことは$f=-1$より不可能。
従って$z=1/(2u)$、$f_y=0$から$y=-1/(8u)$となり、
$f=(1-16u)/(64u^2)=0$を得る。
従って$u\ne1/16$では滑らかであり、
$u=1/16$では唯一の特異点$[1:-2:8]$を持つ。
この点の二次項は非退化（$x=1$でのHessianの行列式は$1/4$）
なので通常節点である。
唯一の通常節点を持つ平面三次曲線は既約被約である：
非被約なら重複成分上で特異となり、
被約可約なら直線と円錐曲線の交叉長は二なので
二つの特異点または普通節点でない一つの特異点を持ち、
三直線の場合にも一つの普通節点だけにはならない。

滑らかな場合はいずれも三線形式が一次独立なので
定理\ref{thm:ou-smooth-determinant-reconstruction}を適用できる。
$u=1/16$の特異点の六座標は
\[
 [2:1/2:-1:1:-2:8]
\]
で全て非零。従って偶奇座標の交点は全て滑らかである。
また各$C_k$は零成分を持ち、零行も零列も持たない。
特異点では$A_{ij}$が全て非零なので、
$\mathscr M_k$にも零成分があり零行・零列はない。
零行・零列を持たないrank一行列には零成分は存在しないので
$\mathscr M_k$はrank一以下にならない。
行列式は零だからrank二である。
従って定理\ref{thm:ou-integral-determinant-criterion}を適用できる。
全ての認識で標準noetherian OU比較は別途仮定する。

台の同値で許した奇置換は巡回の反転を含み得るが、
九線束表示を持つ三ブロックdirected algebraの反対代数も
九線束表示を持つ。実際、列を反転して各線束を双対化し、
全体を次数二線束で捩れば全次数は再び$0,1,2$となる。
この操作はHomと合成を反転させ、直交性・前向きExt消滅も保つ。
従って元の向きの関係代数について曲線表示を得るのであり、
内在的標識を反対圏に移す必要はない。
\end{proof}

\begin{remark}
二つの係数行列、零平面、三次式、特異パラメータと
特異点のHessianおよび六座標は
\path{computations/verify_ou_seven_zero_curves.py}で検算できる。
一般の係数に対する認識は上記の証明による。
$z=7$で未判定の範囲は台116309742の$u=1$と、
台117368558, 125531767, 128957047, 129972855, 133118583の
五つの一係数族である。これは必要条件による残余であり、
全逆Serre反復のvery strong性、標準比較、内在的標識の実在は
主張しない。
\end{remark}

次はこれらの七零係数の残余型に対し、既約特異切断の格子条件、
可約曲線・別foundationによる表示、または高次syzygy・内在的障害を
調べる。さらに一個から六個の零係数型と標準noetherian OU比較の
適用範囲が残る。
